\input{../tex-inputs/latex-leseansicht-vorspann}

\section[Max Schwarzkopf an Arthur Schnitzler, 18. 6. 1904]{L04312 Max Schwarzkopf an Arthur Schnitzler, 18. 6. 1904}
\nopagebreak\mylabel{L04312v}
\rehead{ }\normalsize\beginnumbering\briefempfaengerindex{Schnitzler, Arthur@\textsc{Schnitzler, Arthur}!zzzSchwarzkopf, Max@\emph{von Max Schwarzkopf}!1904-06-181@{18. 6. 1904}|(be}
\toendnotes[C]{\smallbreak\pagebreak[2]}
\correspDesc{Versand  durch Max Schwarzkopf am 18. 6. 1904 in Wien
\newline{}Übermittlung  am 18. 6. 1904 in Wien
\newline{}Zustellung  am 18. 6. 1904 in Wien
\newline{}Erhalt  durch Arthur Schnitzler am 18. 6. 1904 in Wien}\toendnotes[C]{\smallbreak}
\Standort{CUL, Schnitzler, B 96a.}
\physDesc{Postkarte, 211 Zeichen
\newline{}Handschrift: schwarze Tinte, deutsche Kurrent
\newline{}Versand: 1) Stempel: »\nobreak{}\oindex{I., Innere Stadt@\textbf{I., Innere Stadt}, \emph{Verwaltungsgebiet}|pwk}1/1 Wien \textcolor{gray}{9}, 18 6 04, 10–11V\nobreak{}«.   2) Stempel: »\nobreak{}\oindex{I., Innere Stadt@\textbf{I., Innere Stadt}, \emph{Verwaltungsgebiet}|pwk}18/1 Wien 110, 18. 6. 04, 3.N, bestellt\nobreak{}«. }\toendnotes[C]{\smallbreak}\pstart{}{\pb}Herrn
                  \textsc{D\textsuperscript{r} Arthur Schnitzler}\pend{}\pstart{}Wien, XVIII. Währing\oindex{XVIII., Währing@\textbf{XVIII., Währing}, \emph{Verwaltungsgebiet}|pw}.\pend{}\pstart{}Spöttelgaſſe 7\oindex{Wien@\textbf{Wien}!XVIII., Währing@\textbf{XVIII., Währing}!Edmund-Weiß-Gasse 7@\textbf{Edmund-Weiß-Gasse 7}, \emph{Wohngebäude}|pw}\pend{}{\bigskip}\vspace{1em}
\pstart{}{\pb}Sehr geehrter Freund!\pend\vspace{0.5em}
\pstart
           \label{K_L04312-1v}\edtext{der reine Tor\pwindex{Schwarzkopf, Max 12.\,6.\,1857 Wien – 14.\,4.\,1928 ebd.@\textsc{Schwarzkopf, Max} (12.\,6.\,1857 Wien – 14.\,4.\,1928 ebd.), \emph{Rechtsanwalt}!reine Tor. Gesellschaftsstück in vier Akten@\strich\emph{Der reine Tor. Gesellschaftsstück in vier Akten}|pw} hat, wie ich glaube,
               Ihnen keine Schande}{\lemma{\textnormal{\emph{der … Schande}}}\Cendnote{\textnormal{
                  Die Uraufführung von \emph{Der reine Tor}\pwindex{Schwarzkopf, Max 12.\,6.\,1857 Wien – 14.\,4.\,1928 ebd.@\textsc{Schwarzkopf, Max} (12.\,6.\,1857 Wien – 14.\,4.\,1928 ebd.), \emph{Rechtsanwalt}!reine Tor. Gesellschaftsstück in vier Akten@\strich\emph{Der reine Tor. Gesellschaftsstück in vier Akten}|pwk}\eventindex{Neues Deutsches Theater@\textbf{Neues Deutsches Theater}!Uraufführung von Der reine Tor, 17.6.1904@Uraufführung von Der reine Tor, 17.6.1904|pwk} hatte am 17. 6. 1904 am \emph{Neues Deutsches Theater}\orgindex{Neues Deutsches Theater@Neues Deutsches Theater|pwk} in Prag\oindex{Prag@\textbf{Prag}, \emph{Land}|pwk} stattgefunden.
               }}}\label{K_L04312-1} gemacht.\pend
           
\pstart
           Es grüßt Sie und Ihre Frau\pwindex{Schnitzler, Olga 17.\,1.\,1882 Wien – 13.\,1.\,1970 Lugano@\textsc{Schnitzler, Olga} (17.\,1.\,1882 Wien – 13.\,1.\,1970 Lugano), \emph{Schauspielerin, Sängerin}|pwv} vielmals{\\[\baselineskip]}Ihr ergebener{\\[\baselineskip]}\spacefill\mbox{D\textsuperscript{r} M Schwarzkopf}\pend
           \leftskip=0em{}
\pstart
           \raggedleft{}18/VI 04.\pend
           \selectlanguage{ngerman}\endnumbering\briefempfaengerindex{Schnitzler, Arthur@\textsc{Schnitzler, Arthur}!zzzSchwarzkopf, Max@\emph{von Max Schwarzkopf}!1904-06-181@{18. 6. 1904}|)be}\mylabel{L04312h}
\begin{anhang}
\end{anhang}\newcommand{\dateiname}{L04312}\newcommand{\titel}{Max Schwarzkopf an Arthur Schnitzler, 18. 6. 1904}\newcommand{\editorInnen}{Herausgegeben von Jahnke, SelmaMüller, Martin Anton}\input{../tex-inputs/latex-leseansicht-abspann}
